\section{Reference Implementation and Budgets}
\label{app:implementation}

\subsection{Actual Representation-Screen Configuration}
The reference D model uses \texttt{se3wam\_v2/split\_ffn}, four transformer
layers, width 768, eight attention heads, and factorized attention.
DINOv3 ViT-L/16 is frozen \cite{dinov3}. A $224\times224$ agent-view image
produces 196 current visual tokens. H0V has zero visual-history query tokens.
The information track also supplies exact depth and camera geometry.
Robot input retains eight EEF7 states at 20\,Hz. Eight object tokens encode
current GT pose and seven motion intervals using 25 numeric statistics and
pose/interval/acceleration validity flags.

The object statistics include position, rotation vector, interval translation
and relative rotation, elapsed time, linear/angular velocity, and
linear/angular acceleration. Their source is simulator truth in this track.
The current-pose history track removes historical object poses and all derived
interval kinematics while preserving the current pose and robot history.
This is a different information condition, not a different name for the same
input.

D shares attention QKV/output projections and splits feed-forward networks
by world/meta, flow, and action stream. Scene tokens access scene tokens;
meta tokens access scene and meta; flow accesses scene and flow; action
accesses meta and action. Consequently, flow supervision can change shared
weights without an online flow-to-action path. The deterministic pose head
pools sixteen final meta tokens. Calling the model world--action does not
change these actual access relations.

\subsection{Coordinates, Losses, and Generation}
Target future translation is cumulative displacement in the fixed metric
frame. Rotation is the left-relative vector
$\phi_h=\Log(R_{t+\delta_h}R_t^\top)^\vee$.
For valid targets $\mathcal V$, the pose loss is
\begin{equation}
\begin{split}
 \mathcal L_{\rm pose}=\frac{1}{|\mathcal V|}\sum_{h\in\mathcal V}
 \big[&\|\hat d_h-d_h\|_2^2\\
 &+0.1\,d_{\mathrm{SO}(3)}(\hat R_{t+\delta_h},R_{t+\delta_h})^2\big].
\end{split}
\end{equation}
The external weight multiplying this loss is separate from the internal
rotation coefficient. A full twist exponential would couple translation
to rotation through its Jacobian and is not the implemented parameterization.
The reported configuration does not include the former draft's heteroscedastic
pose likelihood, uncertainty bridge, symmetry-minimized loss, or pose--flow
consistency objective.

The action has eight E7 steps: base-frame translation increment, left-relative
rotation vector, and gripper command. Translation and rotation training bounds
are 0.05\,m and 0.5\,rad per axis. The inherited evaluation translation
limit is 0.10\,m per axis. Action flow matching uses a unit Gaussian in
training-normalized action coordinates. Scene flow uses a metric Gaussian
source with standard deviation 0.02\,m. Deterministic auxiliary readouts,
including pose, do not consume flow-matching source noise and time.

Only the action generator runs during Q2 deployment, using five Euler steps.
Conditioning caching is on; temporal/velocity caching is off.
Each replan executes one command. Trainable initialization is fresh
for seed 0 rather than inherited from a trained common-reference checkpoint.
Global batch is 16. Four epochs produce 40,104 updates and 641,664 chunks.
P00, P02, and P05 finish with an eight-device topology; other arms finish
with sixteen. Some retained runs have recorded topology transitions.
These histories preserve the declared exposure and endpoint but do not
establish bitwise equivalence across topology.

\subsection{Alternative Target Definitions}
P04 uses a frozen Wan video VAE target. Nine RGB input frames produce three
temporal latents in the target construction; this is not eight independent
future latent frames. P05 uses dense cumulative 2-D displacement. P06 uses
real codec block motion vectors with reference, direction, and validity.
P07 uses sparse scene tracks and visibility, P08 pooled object DINO features,
and P09 target tracks and visibility. P08/P09 adapt mechanisms rather than
reproduce an entire published model.

Aalpha adds motion-magnitude MSE with weight 0.5 while preserving P03's
Alpha-head initialization. P12, P13, and P14 add future depth, spatial DINO,
and dynamic-region RGB supervision. P15 combines these three at their
individual weights. A non-existent head has an inapplicable readout,
not a zero error. The combination cannot identify every pairwise interaction.
Raw total losses across these objectives are not commensurate.

\subsection{System-Reference Training}
% Transcribed from Research BASE1 and the FUTURE0/BASE2 seed-0 reports.
\begin{table*}[t]
\centering\small
\caption{\tablelang{Training exposure and development counts for all system-reference
endpoints. Batch is global; all endpoints use training seed 0. Initialization
and execution are part of each adapted system, not matched factors.
The two DreamVLA endpoints are independent training schedules.
The 1,260-scene counts appear in the main system table.}{所有系统参照端点的训练量和开发集成功数。Batch 为全局批大小；所有端点的训练 seed 均为 0。初始化与执行方式属于各适配系统自身的设置，并非匹配因素。两个 DreamVLA 端点采用独立训练计划。1,260 场补集的计数见正文系统比较表。}}
\label{tab:system_budgets}
\begin{tabularx}{\textwidth}{@{}lrrrX@{}}
\toprule
\tablelang{System}{系统} & \tablelang{Updates}{更新次数} & \tablelang{Batch}{批大小} & \tablelang{Success /140}{成功数 /140} & \tablelang{Execution / initialization note}{执行与初始化说明} \\
\midrule
ACT & 50,000 & 32 & 50 & \tablelang{Scratch; temporal aggregation}{从头训练；时间聚合} \\
Diffusion Policy & 50,000 & 32 & 54 & \tablelang{Scratch; execute 8 actions}{从头训练；执行 8 个动作} \\
$\pi_{0.5}$ & 30,000 & 256 & 89 & \tablelang{\texttt{pi05\_base}; execute 10}{\texttt{pi05\_base}；执行 10 个动作} \\
DynamicVLA & 50,000 & 44 & 24 & \tablelang{Official DOM checkpoint; native continuous pipeline}{官方 DOM 检查点；原生连续流水线} \\
OpenWAM-Alpha & 34,490 & 192 & 76 & \tablelang{Official Alpha foundation; execute 32}{官方 Alpha 基础模型；执行 32 个动作} \\
DynamicWAM & 40,000 & 128 & 72 & \tablelang{Official full initialization; execute 16}{官方完整初始化；执行 16 个动作} \\
\tablelang{DreamVLA: 20 epochs}{DreamVLA：20 轮} & 48,220 & 512 & 112 & \tablelang{Independent schedule; aggregate forecasts, execute 1}{独立训练计划；聚合预测，执行 1 个动作} \\
\hspace{1em}\tablelang{40 epochs}{40 轮} & 96,440 & 512 & 114 & \tablelang{Independent longer schedule}{独立的更长训练计划} \\
\midrule
\tablelang{Study P00}{本研究 P00} & 40,104 & 16 & 88 & \tablelang{Fresh task network; frozen DINO-L; execute 1}{新训任务网络；冻结 DINO-L；执行 1 个动作} \\
\tablelang{Study P03}{本研究 P03} & 40,104 & 16 & 93 & \tablelang{Same reference recipe as P00}{与 P00 相同的参照配方} \\
\tablelang{Study I}{本研究 I} & 40,104 & 16 & 96 & \tablelang{Fresh CNN condition; full visual history; execute 1}{新训 CNN 条件；完整视觉历史；执行 1 个动作} \\
\bottomrule
\end{tabularx}
\end{table*}

ACT and Diffusion Policy train from scratch for 50,000 updates at global batch
32. The $\pi_{0.5}$ adaptation starts from \texttt{pi05\_base} and trains
30,000 updates at global batch 256 \cite{pi05}. DynamicVLA uses 50,000 updates
at global batch 44. OpenWAM uses a fresh fine-tuning run from the official
Alpha foundation, for five epochs,
34,490 updates at global batch 192, with its declared foundation components.
DynamicWAM uses its official full initialization and 40,000 updates at
global batch 128 \cite{openwam_alpha,dynamicwam}.
DreamVLA's 20-epoch and 40-epoch schedules use independent training runs,
48,220 and 96,440 updates at global batch 512; the shorter endpoint is not
the midpoint checkpoint of the longer schedule \cite{dreamvla}.

Executors also differ: ACT uses temporal aggregation, Diffusion Policy executes
eight actions, $\pi_{0.5}$ ten, and DynamicWAM sixteen before replanning.
The internal representation screen executes one. System-level interpretation
retains these differences. Matched-budget comparisons must also report each
system's inputs, initialization, and execution configuration.
